\begin{spacing}{1.2}

\vspace{0.5em}

When a generative model can draft, code, and summarize on demand, curricula built around information transfer become fragile; \textbf{what remains valuable is the structure of work---tasks, judgment, and responsibility}. We therefore model the Gen-AI transition from the ground up: occupations are decomposed into an auditable \textbf{task DNA} using O*NET (task importance and frequency), and Gen-AI impact is expressed on a \textbf{nine-dimension capability scale} that separates \emph{substitution pressure} from \emph{complementarity gains}. An adoption-driven dynamic then converts these task-level signals into interpretable occupational trajectories (which tasks shrink, which expand, and when the occupational boundary effectively changes).

To ensure our conclusions are not tailored to a single niche, we select three representative occupations by identifying an \textbf{``occupational triangle''} with maximum heterogeneity: \emph{Robotics Engineers} (SOC 17-2199.08), \emph{First-Line Supervisors of Mechanics, Installers, and Repairers} (SOC 49-1011.00), and \emph{Court Reporters and Simultaneous Captioners} (SOC 27-3092.00). \textbf{The selection is stable under candidate-pool expansion} and yields three distinct futures: a robotics pathway dominated by complementarity and task reallocation; a skilled-trade supervisory pathway where routine documentation and scheduling are compressed while field judgment and coordination become more central; and a court-reporting pathway where core production is more directly exposed to substitution, forcing a stronger toolchain redesign.

We translate these occupational insights into institutional policy and curriculum action. For each occupation we recommend one program that can realistically institutionalize Gen-AI resources (MRSD at CMU; UTI Diesel; Generations College Court Reporting) and then evaluate each course using a \textbf{three-axis meta-model}: (i) marginal employability contribution based on O*NET task coverage, (ii) humanistic value spanning civic responsibility, agency, well-being, and lifelong learning, and (iii) AI-overlap v2, defined by each course's nine-dimension footprint weighted by the previously estimated Gen-AI capability strengths. The resulting \textbf{KEEP/TRANSFORM/PRUNE} decisions are designed for portability (course title + description suffices) and emphasize \textbf{transforming high-overlap courses into human--AI collaboration}, evidence-based reasoning, and ethical application rather than indiscriminate deletion.

Finally, we test the most consequential uncertainty---the definition of the target occupation set---using narrow/medium/wide variants. Across all three programs, \textbf{course-level decisions and the set of courses flagged for transformation remain unchanged} (label stability and transformation-set agreement are both 1.00), while fine-grained employability rankings can vary, establishing a clear generalization boundary: \textbf{the framework transfers across institutions, but the intended employment target must be explicitly re-specified when extrapolating}.

\end{spacing}

\begin{keywords}
Generative AI; Task DNA; Employability alignment; Humanistic 
\end{keywords}
